\documentclass[11pt]{article}
\usepackage[margin=1in]{geometry}
\usepackage{amsmath}
\usepackage{amssymb}
\usepackage{graphicx}
\usepackage{booktabs}
\usepackage{hyperref}
\usepackage{biblatex}
\addbibresource{ProjectReferences.bib}

\title{Project 2: Nonlinear ODEs}
\author{Benjamin Cleaver}
\date{\today}
\begin{document}
\maketitle

\textbf{AI policy: }
No AI was used in the creation of this project. Auto-fill suggestions were used for things that were repetitive and I was going to type them anyways. \\


\textbf{Background Information: } Nonlinear ordinary differential equations (ODEs) are used to show changes due to the passage of time.
Some examples of systems that use ODEs are:
\begin{center}
Swinging pendulums, Predator-prey populations, chemical oscillators.
\end{center}

\noindent A system of first-order ODEs will have the form:

\begin{center}
    $\frac{dx}{dt}=f(x,t), \quad x(0)=x_0$
\end{center}
Where x(t) is the state of the system at time t. \\

\noindent In this project, the two systems that I will be demonstrating are the Lorenz System and the Van der Pol oscillator \\

\noindent The Van der Pol oscillator is a Two-Dimension system, while the Lorenz system is a Three-Dimensional system. \\

\section{Lorenz system: } \cite{LorenzSystem}

\subsection{Write the first-order ODEs}
\noindent The Lorenz system can be written as the first-order ODE:

\begin{center}
    $x'= \sigma(y-x), \quad y'=x(\rho-z)-y,\quad z'=xy-\beta z$
\end{center}
\noindent With the solutions $\sigma=10, p=28, and \,\beta=8/3$ \\

\noindent In this system:
\begin{center}
    x represents the rate of convection.\\
    y represents the temperature difference between rising and descending air\\
    z represents the deviation of the temperature from the average.\\
    $\sigma$ represents the fluid viscosity/thermal diffusivity (Prandtl number)\\
    $\rho$ is the Rayleigh number which drives convection\\
    $\beta$ is a system parameter which relates to the geometry of the system.
\end{center}

\noindent Essentially, it describes the behavior of convection currents in the atmosphere. \\

\subsection{Write the right-hand side and solve with solve\_ivp}

\noindent Using this python code:
\begin{center}
    s=10.0 ; p=28.0 ; b=8.0/3.0 \\
sol=solve\_ivp(f,(0,50), [.01,.01,.01],t\_eval=tt)
\end{center}
\noindent This gives us the right-hand side (f) on an interval (0,50) in the provided list of starting values [.01,.01,.01]

\subsection{Plot solutions against time for a few initial states and parameter values}

\begin{center}
    fig,ax=plt.subplots(figsize=(10,6)) \\
plt.axhline(x2, color='black',lw=1,linestyle='--',label=f"Steady state: x={x2:.2f}")
plt.axhline(-x2, color='black',lw=1,linestyle='--',label=f"Steady state: x={-x2:.2f}") \\

ax.plot(sol.t,sol.y[0],'-',color='purple',lw=1,label='Trajectory') \\

ax.set\_xlabel('t'); ax.set\_ylabel('x'); ax.set\_title('Lorenz System') \\
ax.legend() \\
plt.show()
\end{center}

\begin{figure}[h]
    \centering
    \includegraphics[width=1\linewidth]{LorenzSystem.png}
    \caption{Lorenz System}
    \label{fig:placeholder}
\end{figure}

\subsection{Plot error at a fixed final time against the tolerance rtol on a log-log scale}

\begin{center}
    print("rtol      max error   calls of f") \\
for rtol in [1e-3, 1e-6, 1e-9]: \\
    sol = solve\_ivp(f, (0, 50), [.01,.01,.01], t\_eval=tt, rtol=rtol, atol=1e-12) \\
    err = np.max(np.abs(sol.y[0] - y\_exact(tt))) \\
    print(f"{rtol:8.0e}   {err:9.1e}   {sol.nfev:5d}")
\end{center}

\noindent When $|y|$ is large, the rtol controls the accuracy; when $|y|$ is small, atol does. A smaller tolerance results in smaller errors.
\begin{center}
\noindent rtol    \quad  max error  \quad calls of f \\
1e-03  \quad   1.9e+01  \quad  4634 \\
1e-06  \quad   1.9e+01  \quad 13520 \\
1e-09  \quad   1.9e+01 \quad  49100
\end{center}

\section{Van der Pol oscillator} \cite{VanDerPol}

\subsection{Write the first-order ODEs}
\noindent The Van der Pol Oscillator system can be written as the first-order ODE:

\begin{center}
    tt=np.linspace(0,50,400)\\
u=1.0\\
def f(t, xy): \\
    x,y=xy \\
    return [y, u*(1-x**2)*y-x] \\
\end{center}

In this system:
\begin{center}
    $\mu$ represents the dampening.
\end{center}

\subsection{Write the right-hand side and solve with solve\_ivp}

Using this python code:

\begin{center}
    sol = solve\_ivp(f,(0,50),[1,0], t\_eval=tt)
\end{center}

We can find the right-hand side (f) on an interval (0,50) in the provided list of starting values [1,0]

\subsection{Plot solutions against time for a few initial states and parameter values}

\begin{center}
    fig,ax=plt.subplots(figsize=(10,6))\\

ax.plot(sol.t, sol.y[0], color='purple', lw=1, label='x(t)')\\

ax.set\_xlabel('t')\\
ax.set\_ylabel('x')\\
ax.set\_title('Van der Pol Oscillator')\\
ax.legend()\\

plt.show()\\
\end{center}

\begin{figure}[h]
    \centering
    \includegraphics[width=1\linewidth]{VanDerPol.png}
    \caption{Van der Pol Oscillator}
    \label{fig:placeholder}
\end{figure}

\subsection{Plot error at a fixed final time against the tolerance rtol on a log-log scale}

\begin{center}
    print("   rtol      max error   calls of f")
for rtol in [1e-3, 1e-6, 1e-9]:\\
    sol\_exact = solve\_ivp(f, (0, 50), [1,0], t\_eval=tt, rtol=rtol, atol=1e-12)\\
    err = np.max(np.abs(sol.y[0] - sol\_exact.y[0]))\\
    print(f"{rtol:8.0e}   {err:9.1e}   {sol.nfev:5d}")
\end{center}

\noindent When $|y|$ is large, the rtol controls the accuracy; when $|y|$ is small, atol does. A smaller tolerance results in smaller errors.

\begin{center}
rtol  \quad    max error  \quad calls of f \\
   1e-03   \quad  2.2e-02  \quad  1106 \\
   1e-06 \quad 2.7e-02  \quad  1106 \\
   1e-09  \quad  2.7e-02  \quad  1106
\end{center}

\section{Conclusion:}

\subsection{References:}

\printbibliography



\end{document}